\subsection{Orthogonal Complements}

% **Role:** linear algebra with metric and angle structure.
%
% Orthogonal complements, projection theorem, orthogonal projections, best approximation.

Throughout this subsection, $\mathbb{F}$ is $\mathbb{R}$ or $\mathbb{C}$ and $\mathcal{V}$ is an inner-product space over $\mathbb{F}$.

Orthogonal complements turn orthogonality into a subspace operation and organize canonical direct-sum decompositions.

\begin{definition} \label{def:orthogonal-subsets-and-direct-sums} \hlt{Orthogonality and Orthogonal Direct Sums}\\
Let $\mathcal{X},\mathcal{Y}\subseteq\mathcal{V}$ and let $\mathcal{S},\mathcal{T}\leq\mathcal{V}$.
\begin{enumerate}[label=(\roman*)]
\setlength{\itemsep}{0pt}
\item $\mathcal{X}$ and $\mathcal{Y}$ are \hlt{orthogonal}, written $\mathcal{X}\perp\mathcal{Y}$, if $(\forall\,\mathbf{x}\in\mathcal{X})(\forall\,\mathbf{y}\in\mathcal{Y})\,\innerprod{\mathbf{x},\mathbf{y}}=0$.
\item $\mathcal{V}$ is the \hlt{orthogonal direct sum} of $\mathcal{S}$ and $\mathcal{T}$, written $\mathcal{V}=\mathcal{S}\odot\mathcal{T}$, if $\mathcal{V}=\mathcal{S}\oplus\mathcal{T}$ and $\mathcal{S}\perp\mathcal{T}$.
\end{enumerate}
\end{definition}

\begin{definition} \label{def:orthogonal-complement} \hlt{Orthogonal Complement}\\
For $\mathcal{U}\subseteq\mathcal{V}$, its \hlt{orthogonal complement} is $\mathcal{U}^{\perp}\defeq\{\mathbf{v}\in\mathcal{V}:(\forall\,\mathbf{u}\in\mathcal{U})\,\innerprod{\mathbf{v},\mathbf{u}}=0\}$.
\end{definition}

\begin{proposition} \label{prop:basic-properties-orthogonal-complement} \hlt{Basic Properties of Orthogonal Complements}\\
For subsets $\mathcal{U},\mathcal{W}\subseteq\mathcal{V}$,
\begin{enumerate}[label=(\roman*)]
\setlength{\itemsep}{0pt}
\item $\mathcal{U}^{\perp}$ is a subspace of $\mathcal{V}$;
\item $\{\mathbf{0}\}^{\perp}=\mathcal{V}$;
\item $\mathcal{V}^{\perp}=\{\mathbf{0}\}$;
\item $\mathcal{U}\cap\mathcal{U}^{\perp}\subseteq\{\mathbf{0}\}$;
\item if $\mathcal{U}\subseteq\mathcal{W}$, then $\mathcal{W}^{\perp}\subseteq\mathcal{U}^{\perp}$.
\end{enumerate}
\begin{proof}
\begin{enumerate}[label=(\roman*)]
\setlength{\itemsep}{0pt}
\item $\mathbf{0}\in\mathcal{U}^{\perp}$ by Proposition \ref{prop:basic-inner-product-identities}(ii). Fix $\mathbf{v},\mathbf{w}\in\mathcal{U}^{\perp}$ and $a\in\mathbb{F}$.\\
Then $(\forall\,\mathbf{u}\in\mathcal{U})\,\innerprod{a\mathbf{v}+\mathbf{w},\mathbf{u}}=a\innerprod{\mathbf{v},\mathbf{u}}+\innerprod{\mathbf{w},\mathbf{u}}=0$ by linearity, so $a\mathbf{v}+\mathbf{w}\in\mathcal{U}^{\perp}$. Hence $\mathcal{U}^{\perp}\leq\mathcal{V}$ by Theorem \ref{thm:subspacetest}.
\item $(\forall\,\mathbf{v}\in\mathcal{V})\,\innerprod{\mathbf{v},\mathbf{0}}=0$ by Proposition \ref{prop:basic-inner-product-identities}(ii), so $\{\mathbf{0}\}^{\perp}=\mathcal{V}$.
\item $\mathbf{0}\in\mathcal{V}^{\perp}$ by (i). If $\mathbf{v}\in\mathcal{V}^{\perp}$, then $\norm{\mathbf{v}}^2=\innerprod{\mathbf{v},\mathbf{v}}=0$, so $\mathbf{v}=\mathbf{0}$ by definiteness.
\item If $\mathbf{v}\in\mathcal{U}\cap\mathcal{U}^{\perp}$, then $\norm{\mathbf{v}}^2=\innerprod{\mathbf{v},\mathbf{v}}=0$, so $\mathbf{v}=\mathbf{0}$ by definiteness.
\item Suppose $\mathcal{U}\subseteq\mathcal{W}$ and fix $\mathbf{v}\in\mathcal{W}^{\perp}$. Then $(\forall\,\mathbf{u}\in\mathcal{U})\,\innerprod{\mathbf{v},\mathbf{u}}=0$, so $\mathbf{v}\in\mathcal{U}^{\perp}$.
\end{enumerate}
\hlt{Structural Commentary}: Subspace closure, definiteness, and reversed constraint inclusion establish the five properties.\\
\hlt{Key Insight}: Larger tested sets impose stronger orthogonality constraints.
\end{proof}
\end{proposition}

\begin{remark} \label{rem:orthogonal-complements-definiteness-warning} \hlt{Definiteness Matters}\\
For a bilinear form $B:\mathcal{V}\times\mathcal{V}\to\mathbb{F}$ and $\mathcal{U}\subseteq\mathcal{V}$, set $\mathcal{U}^{\perp_B}\defeq\{\mathbf{v}\in\mathcal{V}:(\forall\,\mathbf{u}\in\mathcal{U})\,B(\mathbf{v},\mathbf{u})=0\}$.\\
Bilinearity makes $\mathcal{U}^{\perp_B}$ a subspace, but without definiteness one may have $\mathcal{V}^{\perp_B}\neq\{\mathbf{0}\}$ or $\mathcal{U}\cap\mathcal{U}^{\perp_B}\neq\{\mathbf{0}\}$.\\
Thus $\mathcal{V}=\mathcal{U}\oplus\mathcal{U}^{\perp_B}$ is an inner-product conclusion, not a consequence of bilinearity alone.
\end{remark}

\begin{theorem} \label{thm:orthogonal-direct-sum-decomposition} \hlt{Direct Sum with the Orthogonal Complement}\\
If $\mathcal{U}$ is a finite-dimensional subspace of $\mathcal{V}$, then $\mathcal{V}=\mathcal{U}\oplus\mathcal{U}^{\perp}$.
\begin{proof}
If $\mathcal{U}=\{\mathbf{0}\}$, then $\mathcal{U}^{\perp}=\mathcal{V}$ by Proposition \ref{prop:basic-properties-orthogonal-complement}(ii), and hence $\mathcal{V}=\mathcal{U}\oplus\mathcal{U}^{\perp}$.\\
Suppose $\dim(\mathcal{U})=m\geq1$ and choose an orthonormal basis $(\mathbf{e}_1,\ldots,\mathbf{e}_m)$ by Corollary \ref{cor:finite-dimensional-inner-product-space-has-orthonormal-basis}.\\
Fix $\mathbf{v}\in\mathcal{V}$ and set $\mathbf{u}\defeq\sum\limits_{j=1}^{m}\innerprod{\mathbf{v},\mathbf{e}_j}\mathbf{e}_j\in\mathcal{U}$ and $\mathbf{w}\defeq\mathbf{v}-\mathbf{u}$.\\
For $k\in\{1,\ldots,m\}$, $\innerprod{\mathbf{w},\mathbf{e}_k}=\innerprod{\mathbf{v},\mathbf{e}_k}-\sum\limits_{j=1}^{m}\innerprod{\mathbf{v},\mathbf{e}_j}\innerprod{\mathbf{e}_j,\mathbf{e}_k}=0$ by orthonormality.\\
Fix $\mathbf{x}\in\mathcal{U}$ and write $\mathbf{x}=\sum\limits_{k=1}^{m}a_k\mathbf{e}_k$. Then $\innerprod{\mathbf{w},\mathbf{x}}=\sum\limits_{k=1}^{m}\overline{a_k}\innerprod{\mathbf{w},\mathbf{e}_k}=0$ by Proposition \ref{prop:basic-inner-product-identities}(iii)--(iv), so $\mathbf{w}\in\mathcal{U}^{\perp}$ and $\mathbf{v}=\mathbf{u}+\mathbf{w}$.\\
Hence $\mathcal{V}=\mathcal{U}+\mathcal{U}^{\perp}$. The identity $\mathcal{U}\cap\mathcal{U}^{\perp}=\{\mathbf{0}\}$ holds by Proposition \ref{prop:basic-properties-orthogonal-complement}(i),(iv), so the sum is direct.\\
\hlt{Structural Commentary}: An orthonormal expansion constructs the component; definiteness makes the decomposition direct.\\
\hlt{Key Insight}: Every vector splits into a subspace component and orthogonal residual.
\end{proof}
\end{theorem}

\begin{proposition} \label{prop:orthogonal-direct-sum-characterization} \hlt{Orthogonal Direct Sum Characterization}\\
For $\mathcal{S},\mathcal{T}\leq\mathcal{V}$, the following are equivalent:
\begin{enumerate}[label=(\roman*)]
\setlength{\itemsep}{0pt}
\item $\mathcal{V}=\mathcal{S}\odot\mathcal{T}$.
\item $\mathcal{V}=\mathcal{S}\oplus\mathcal{T}$ and $\mathcal{T}=\mathcal{S}^{\perp}$.
\end{enumerate}
\begin{proof}
\begin{enumerate}[label=(\roman*)]
\setlength{\itemsep}{0pt}
\item[$(i)$]$\Rightarrow (ii)$ Suppose $\mathcal{V}=\mathcal{S}\odot\mathcal{T}$. Then $\mathcal{V}=\mathcal{S}\oplus\mathcal{T}$ by orthogonal-direct-sum definition, and $\mathcal{T}\subseteq\mathcal{S}^{\perp}$ by Proposition \ref{prop:orthogonality-and-zero}(i).\\
Fix $\mathbf{v}\in\mathcal{S}^{\perp}$ and write $\mathbf{v}=\mathbf{s}+\mathbf{t}$ with $\mathbf{s}\in\mathcal{S}$ and $\mathbf{t}\in\mathcal{T}$.\\
Then $\mathbf{s}=\mathbf{v}-\mathbf{t}\in\mathcal{S}^{\perp}$, so $\mathbf{s}=\mathbf{0}$ by Proposition \ref{prop:basic-properties-orthogonal-complement}(iv). Hence $\mathbf{v}=\mathbf{t}\in\mathcal{T}$ and $\mathcal{T}=\mathcal{S}^{\perp}$. \item[$(ii)$]$\Rightarrow (i)$ Suppose $\mathcal{V}=\mathcal{S}\oplus\mathcal{T}$ and $\mathcal{T}=\mathcal{S}^{\perp}$. Then $\mathcal{S}\perp\mathcal{T}$ by Proposition \ref{prop:orthogonality-and-zero}(i), so $\mathcal{V}=\mathcal{S}\odot\mathcal{T}$ by orthogonal-direct-sum definition.
\end{enumerate}
\hlt{Structural Commentary}: Direct-sum uniqueness forces an orthogonal partner to be maximal.\\
\hlt{Key Insight}: A full orthogonal summand is exactly the complement.
\end{proof}
\end{proposition}

\begin{corollary} \label{cor:dimension-of-orthogonal-complement} \hlt{Dimension of the Orthogonal Complement}\\
If $\mathcal{V}$ is finite-dimensional and $\mathcal{U}$ is a subspace of $\mathcal{V}$, then $\dim(\mathcal{U}^{\perp})=\dim(\mathcal{V})-\dim(\mathcal{U})$.
\begin{proof}
The subspace $\mathcal{U}$ is finite-dimensional by Corollary \ref{cor:finite-dimensional-spaces-and-subspaces-have-bases}.\\
By Theorem \ref{thm:orthogonal-direct-sum-decomposition}, $\mathcal{V}=\mathcal{U}\oplus\mathcal{U}^{\perp}$, so $\dim(\mathcal{V})=\dim(\mathcal{U})+\dim(\mathcal{U}^{\perp})$ by Corollary \ref{cor:dimension-direct-sum}. Rearranging proves the formula.\\
\hlt{Structural Commentary}: The direct-sum dimension formula converts decomposition into subtraction.\\
\hlt{Key Insight}: The complement contains precisely the unused dimensions.
\end{proof}
\end{corollary}

\begin{corollary} \label{cor:double-orthogonal-complement} \hlt{Double Orthogonal Complement}\\
If $\mathcal{U}$ is a finite-dimensional subspace of $\mathcal{V}$, then $\mathcal{U}=(\mathcal{U}^{\perp})^{\perp}$.
\begin{proof}
Fix $\mathbf{u}\in\mathcal{U}$. For $\mathbf{w}\in\mathcal{U}^{\perp}$, $\innerprod{\mathbf{w},\mathbf{u}}=0$ by membership, so $\innerprod{\mathbf{u},\mathbf{w}}=0$ by Proposition \ref{prop:orthogonality-and-zero}(i). Hence $\mathcal{U}\subseteq(\mathcal{U}^{\perp})^{\perp}$.\\
Conversely, fix $\mathbf{v}\in(\mathcal{U}^{\perp})^{\perp}$. By Theorem \ref{thm:orthogonal-direct-sum-decomposition}, write $\mathbf{v}=\mathbf{u}_0+\mathbf{w}$ with $\mathbf{u}_0\in\mathcal{U}$ and $\mathbf{w}\in\mathcal{U}^{\perp}$.\\
Since $(\mathcal{U}^{\perp})^{\perp}$ is a subspace by Proposition \ref{prop:basic-properties-orthogonal-complement}(i), the established inclusion implies $\mathbf{w}=\mathbf{v}-\mathbf{u}_0\in(\mathcal{U}^{\perp})^{\perp}$.\\
Thus $\mathbf{w}\in\mathcal{U}^{\perp}\cap(\mathcal{U}^{\perp})^{\perp}$, so $\mathbf{w}=\mathbf{0}$ by Proposition \ref{prop:basic-properties-orthogonal-complement}(iv), and $\mathbf{v}=\mathbf{u}_0\in\mathcal{U}$.\\
\hlt{Structural Commentary}: Symmetry establishes one inclusion; the orthogonal decomposition eliminates every excess vector.\\
\hlt{Key Insight}: Double orthogonality adds no new finite-dimensional direction.
\end{proof}
\end{corollary}

\begin{corollary} \label{cor:double-orthogonal-complement-of-finite-span} \hlt{Double Complement of a Finite Span}\\
If $\mathcal{X}\subseteq\mathcal{V}$ and $\operatorname{span}(\mathcal{X})$ is finite-dimensional, then $(\mathcal{X}^{\perp})^{\perp}=\operatorname{span}(\mathcal{X})$.
\begin{proof}
Let $\mathcal{W}=\operatorname{span}(\mathcal{X})$. Since $\mathcal{X}\subseteq\mathcal{W}$, $\mathcal{W}^{\perp}\subseteq\mathcal{X}^{\perp}$ by Proposition \ref{prop:basic-properties-orthogonal-complement}(v).\\
Conversely, fix $\mathbf{v}\in\mathcal{X}^{\perp}$ and $\mathbf{w}\in\mathcal{W}$. Write $\mathbf{w}=\sum\limits_{j=1}^{r}a_j\mathbf{x}_j$ with $\mathbf{x}_j\in\mathcal{X}$.\\
Then $\innerprod{\mathbf{v},\mathbf{w}}=\sum\limits_{j=1}^{r}\overline{a_j}\innerprod{\mathbf{v},\mathbf{x}_j}=0$ by Proposition \ref{prop:basic-inner-product-identities}(iii)--(iv), so $\mathbf{v}\in\mathcal{W}^{\perp}$.\\
Therefore $\mathcal{X}^{\perp}=\mathcal{W}^{\perp}$, and $(\mathcal{X}^{\perp})^{\perp}=(\mathcal{W}^{\perp})^{\perp}=\mathcal{W}$ by Corollary \ref{cor:double-orthogonal-complement}.\\
\hlt{Structural Commentary}: Span invariance reduces the double complement to the preceding corollary.\\
\hlt{Key Insight}: Orthogonality depends on the span of the tested set.
\end{proof}
\end{corollary}

\begin{perspective} \label{per:orthogonal-subspaces-versus-complements} \hlt{Orthogonal Versus Complementary}\\
Let $\mathcal{U},\mathcal{W}\leq\mathcal{V}$ with $\mathcal{U}\perp\mathcal{W}$. By Definitions \ref{def:orthogonal-subsets-and-direct-sums} and \ref{def:orthogonal-complement}, $\mathcal{W}\subseteq\mathcal{U}^{\perp}$ and $\mathcal{U}\subseteq\mathcal{W}^{\perp}$, not necessarily equality.\\
If $\mathcal{V}$ is finite-dimensional and $\dim(\mathcal{U})+\dim(\mathcal{W})=\dim(\mathcal{V})$, then $\dim(\mathcal{U}^{\perp})=\dim(\mathcal{W})$ and $\dim(\mathcal{W}^{\perp})=\dim(\mathcal{U})$ by Corollary \ref{cor:dimension-of-orthogonal-complement}.\\
The two inclusions are therefore equalities by Proposition \ref{prop:dimension-bounds-and-basis-criteria}(ii).
\end{perspective}

Figure \ref{fig:orthogonal-subspaces-vs-complements} contrasts a proper inclusion $\mathcal{W}\subsetneq\mathcal{U}^{\perp}$ with the saturated decomposition $\mathcal{V}=\mathcal{U}\oplus\mathcal{U}^{\perp}$.

\begin{figure}[H]
  \centering
  \begin{tikzpicture}[
  x=0.88cm,
  y=0.88cm,
  >=Stealth,
  every node/.style={font=\small},
  host/.style={draw=gray!50, fill=gray!3, rounded corners=3pt, line width=0.55pt,
    minimum width=4cm, minimum height=2.2cm},
  sub/.style={draw=blue!65!black, fill=blue!7, rounded corners=2pt,
    minimum width=1.2cm, minimum height=0.65cm, align=center},
  orth/.style={draw=teal!65!black, fill=teal!8, rounded corners=2pt,
    minimum width=1.2cm, minimum height=0.65cm, align=center},
  gap/.style={draw=gray!55, fill=gray!8, rounded corners=2pt,
    minimum width=1.2cm, minimum height=0.52cm, align=center},
  shell/.style={draw=teal!55!black, dashed, rounded corners=2pt, line width=0.65pt},
  orthmark/.style={gray!65!black, line width=0.65pt},
  title/.style={font=\small}
]
  \node[host] (left) at (0,0) {};
  \node[anchor=north west] at (-2.05,1.05) {$\mathcal{V}$};
  \node[sub] (u1) at (-1.05,-0.48) {$\mathcal{U}$};
  \node[orth] (w1) at (0.95,-0.48) {$\mathcal{W}$};
  \node[gap] (r1) at (0.95,0.47) {unused};
  \draw[shell] (0.12,-0.98) rectangle (1.78,0.82);
  \node[teal!55!black] at (1.44,0.77) {$\mathcal{U}^{\perp}$};
  \draw[orthmark] (-0.16,-0.08) -- (-0.02,-0.08) -- (-0.02,0.06);
  \node at (0,-1.43) {$\mathcal{W}\subset\mathcal{U}^{\perp}$};

  \node[host] (right) at (6.0,0) {};
  \node[anchor=north west] at (3.95,1.05) {$\mathcal{V}$};
  \node[sub, minimum width=1.52cm, minimum height=1.6cm] (u2) at (5.14,0) {$\mathcal{U}$};
  \node[orth, minimum width=1.52cm, minimum height=1.6cm] (p2) at (6.86,0) {$\mathcal{U}^{\perp}$};
  \draw[orthmark] (5.92,0.12) -- (6.08,0.12) -- (6.08,-0.04);
  \node at (6,-1.43) {$\mathcal{V}=\mathcal{U}\oplus\mathcal{U}^{\perp}$};

  \node[title] at (0,1.48) {orthogonal only};
  \node[title] at (6,1.48) {orthogonal complement};
\end{tikzpicture}

  \caption{Orthogonal subspaces are complements only when their sum fills $\mathcal{V}$.}
  \label{fig:orthogonal-subspaces-vs-complements}
\end{figure}
